\section{A symmetric chiral determinant with genuine motion}
\label{sec:chiral}

The one-sided gauge theorem requires an endpoint-symmetric alternative.
Define the holomorphic operator pencil
\[
 \cB_s=
 \begin{pmatrix}
 0&L_s\\ L_{1-s}^{\mathsf T}&0
 \end{pmatrix}.                              \tag{5.1}
\]
The transpose is taken in the entropy basis; it is linear rather than
antilinear in $s$.

\begin{proposition}[Schatten strip and exact symmetries]
\label{prop:strip}
One has $L_s\in S_q$ whenever $q\RePart s>1$.  Hence
$\cB_s\in S_3$ on
\[
 \frac13<\RePart s<\frac23,
\]
where
\[
 \Delta_3(s,z):=\det{}_3(I-z\cB_s)            \tag{5.2}
\]
is holomorphic.  It satisfies
\[
 \Delta_3(1-s,z)=\Delta_3(s,z),
 \qquad
 \Delta_3(\bar s,\bar z)=\overline{\Delta_3(s,z)}. \tag{5.3}
\]
For $s=1/2+it$, $\cB_s$ is self-adjoint.
\end{proposition}

\begin{proof}
The diagonal and weighted-shift singular values are controlled by the
$\ell^q$ sequence $(p_n^{-\RePart s})_n$.  Both off-diagonal blocks are in
$S_3$ precisely in the displayed common strip.  Furthermore
$\cB_{1-s}=\cB_s^{\mathsf T}$, and regularized determinants are invariant
under transpose.  Real coefficients give conjugation symmetry.  On the
critical line $L_{1-s}^{\mathsf T}=L_s^*$.
\end{proof}

The regularization deletes powers one and two.  Odd chiral traces vanish, so
the first retained coefficient is the fourth trace.  Unlike the Paper06
block-preserving pair, it moves.

Put $u_n=p_n^{-1}$ and, on $s=1/2+it$, define
\[
 x_n(t)=\frac14\left(u_n+u_{n+1}
 +2\sqrt{u_nu_{n+1}}
 \cos\left(t\log\frac{p_{n+1}}{p_n}\right)\right). \tag{5.4}
\]
This is the squared modulus of the $n$th successor coefficient of $L_s$.

\begin{theorem}[Strict fourth-Schatten spectral motion]
\label{thm:motion}
For the infinite SD-C09 transfer,
\[
 \|L_{1/2+it}\|_4^4
 =\sum_n(u_n+x_n(t))^2+2\sum_nu_{n+1}x_n(t), \tag{5.5}
\]
and
\[
 \|L_{1/2+it}\|_4^4<\|L_{1/2}\|_4^4
 \qquad(t\ne0).                              \tag{5.6}
\]
Consequently $\cB_{1/2+it}$ is not unitarily gauge-equivalent to a constant
family and
\[
 \Tr\cB_{1/2+it}^{,4}=2\|L_{1/2+it}\|_4^4
\]
varies with $t$.
\end{theorem}

\begin{proof}
For a lower bidiagonal $L$, the diagonal of $L^*L$ is $u_n+x_n$ and its
nearest off-diagonal squared modulus is $u_{n+1}x_n$.  Squaring and tracing
gives \cref{eq:s4}:
\begin{equation}
 \Tr(L^*L)^2=\sum_n(u_n+x_n)^2+2\sum_nu_{n+1}x_n.
 \label{eq:s4}
\end{equation}
Every summand is strictly increasing in its $x_n\ge0$, and
$x_n(t)\le x_n(0)$.  Equality would require
$t\log(p_{n+1}/p_n)\in2\pi\ZZ$ for every $n$.  In particular, for the
successive ratios $3/2$ and $5/3$, there would be integers $k,\ell$ with
\[
 (3/2)^\ell=(5/3)^k.
\]
Unique factorization forces $k=\ell=0$, and then $t=0$.  Thus at least one
$x_n$ is strictly smaller for every nonzero $t$, proving the claim.
\end{proof}

Near $z=0$, the $\det_3$ expansion begins with
\[
 \log\Delta_3(1/2+it,z)
 =-\frac{z^4}{4}\Tr\cB_{1/2+it}^{,4}+O(z^6)
 =-\frac{z^4}{2}\|L_{1/2+it}\|_4^4+O(z^6). \tag{5.7}
\]
Thus the moving spectrum is visible to the chiral regularized determinant.
It is not visible to the Euler Fredholm determinant of $L_s$.  The two
determinants have overlapping symbolic ancestry but are not an analytic
continuation of one another.
